\documentclass[11pt,a4paper]{article}

\usepackage[T1]{fontenc}
\usepackage[utf8]{inputenc}
\usepackage[brazilian]{babel}
\usepackage{lmodern}
\usepackage[a4paper,margin=2.15cm,headheight=15pt]{geometry}
\usepackage{microtype}
\usepackage{amsmath,amssymb}
\usepackage{graphicx}
\usepackage{booktabs}
\usepackage{tabularx}
\usepackage{longtable}
\usepackage{array}
\usepackage{enumitem}
\usepackage{xcolor}
\usepackage{float}
\usepackage{caption}
\usepackage{fancyhdr}
\usepackage{hyperref}
\usepackage{url}

\definecolor{navy}{HTML}{17365D}
\definecolor{bluebox}{HTML}{EAF2F8}
\definecolor{redbox}{HTML}{FCE8E6}
\definecolor{greenbox}{HTML}{EAF4EA}
\definecolor{amberbox}{HTML}{FFF4D6}
\definecolor{muted}{HTML}{5B6573}

\hypersetup{
  colorlinks=true,
  linkcolor=navy,
  urlcolor=navy,
  citecolor=navy,
  pdftitle={Diagnóstico de generalização do pipeline GLSTM},
  pdfauthor={Auditoria técnica do repositório LSTM-GNN}
}

\pagestyle{fancy}
\fancyhf{}
\lhead{LSTM-GNN --- diagnóstico de generalização}
\rhead{\thepage}
\renewcommand{\headrulewidth}{0.4pt}

\setlist[itemize]{leftmargin=1.45em,itemsep=0.35em,topsep=0.4em}
\setlist[enumerate]{leftmargin=1.65em,itemsep=0.45em,topsep=0.4em}
\setlength{\parindent}{0pt}
\setlength{\parskip}{0.55em}
\setlength{\emergencystretch}{3em}
\captionsetup{font=small,labelfont=bf}

\newcolumntype{Y}{>{\raggedright\arraybackslash}X}
\newcolumntype{C}{>{\centering\arraybackslash}X}
\newcommand{\code}[1]{\texttt{\detokenize{#1}}}
\newcommand{\callout}[2][bluebox]{%
  \par\medskip\noindent
  \colorbox{#1}{\parbox{\dimexpr\linewidth-2\fboxsep\relax}{#2}}%
  \par\medskip
}
\newcommand{\severity}[1]{\textbf{\textcolor{navy}{#1}}}

\begin{document}
\pagenumbering{roman}
\hypersetup{pageanchor=false}

\begin{titlepage}
  \centering
  \vspace*{1.7cm}
  {\Huge\bfseries\color{navy} Diagnóstico de generalização do pipeline GLSTM\par}
  \vspace{0.5cm}
  {\Large Por que a previsão converge para a média --- ou para uma média ponderada\par}
  \vspace{1.2cm}
  \rule{0.84\textwidth}{0.8pt}
  \vspace{0.8cm}

  \begin{minipage}{0.87\textwidth}
  \large
  Auditoria do caminho realmente executado por \code{src/run_experiment.py}, incluindo dados,
  cronologia, janelas, normalização, grafo, arquitetura GLSTM, função objetivo,
  treinamento, checkpoint e inferência.

  \vspace{0.7cm}
  \normalsize
  \textbf{Configuração de referência:}
  \path{Experiments/run_experiment/glstm_sweep_20260715_120845}

  \textbf{Data da auditoria:} 15 de julho de 2026

  \textbf{Escopo:} diagnóstico e recomendações. Nenhuma alteração funcional foi aplicada ao
  código de treinamento.
  \end{minipage}

  \vfill
  {\large Repositório \code{LSTM-GNN}\par}
\end{titlepage}
\hypersetup{pageanchor=true}

\tableofcontents
\clearpage
\pagenumbering{arabic}

\section{Conclusão executiva}

\callout[redbox]{
\textbf{Diagnóstico principal.} O comportamento observado não é, no estado atual, um caso
clássico de rede que ajusta muito bem o treino e falha apenas em dados novos. O $R^2$ ordinário
já é negativo no treino. A rede aprende modestamente a função objetivo configurada, mas essa
função objetivo, combinada com pouco sinal preditivo para D+2--D+5, favorece justamente uma
previsão quase constante. Com MSE comum, essa constante é a média. Com a
\code{quantile_mse} atual, ela é uma \textbf{média ponderada alta}, aproximadamente
$15{,}29$ mm.
}

Os dados e o alinhamento temporal foram verificados de ponta a ponta. Não foram encontrados
NaNs, infinitos, datas faltantes, troca de eixos ou vazamento de alvos entre treino, validação e
teste. O contrato efetivo é

\[
X\in\mathbb{R}^{B\times 35\times 62\times 25}
\quad\longrightarrow\quad
\widehat{y}\in\mathbb{R}^{B\times 5\times 62},
\]

compatível com $y\in\mathbb{R}^{B\times 5\times 62}$. Os scalers são ajustados somente no
treino. Assim, o colapso não é explicado por um erro básico de shape ou cronologia.

As causas mais importantes, em ordem, são:

\begin{enumerate}
  \item \severity{Crítica --- objetivo desalinhado.} A função chamada
  \code{QuantileMSELoss} não é uma quantile loss do tipo pinball. Ela é uma MSE por faixas do
  alvo. Com Q90 e pesos automáticos, cada observação extrema recebe aproximadamente nove vezes
  o peso de uma observação comum. O ótimo sem sinal deixa de ser $4{,}58$ mm e passa a ser
  $15{,}29$ mm.

  \item \severity{Crítica --- densificação do grafo aprendível.} Com
  \code{LOCK_TOPOLOGY=False}, todos os pares de estações tornam-se parametrizados. O mesmo
  \code{weight_decay=0.005} é aplicado aos logits da adjacência; levar um logit negativo para
  zero significa levar a probabilidade de uma não-aresta para $0{,}5$. O checkpoint terminou
  com 69,0\% da massa normalizada em ligações que não existiam no KNN inicial.

  \item \severity{Alta --- informação insuficiente para horizontes longos.} A entrada contém
  somente agregados meteorológicos \emph{passados}. Não há NWP futuro alinhado a cada lead,
  calendário, altitude ou embedding de estação; vento está desligado. D+1 ainda contém algum
  sinal, mas a variabilidade prevista cai rapidamente até D+5.

  \item \severity{Alta --- compressão e suavização.} O mesmo estado oculto final de dimensão
  32 produz os cinco leads; a mistura espacial é aplicada dentro da recorrência a cada um dos
  35 passos. Isso favorece \emph{oversmoothing} temporal e espacial.

  \item \severity{Média --- avaliação incompleta.} Há somente um corte temporal, nenhuma
  validação rolling-origin, nenhum holdout espacial e nenhum baseline climatológico integrado ao
  runner. Além disso, alguns arquivos misturam métricas padronizadas e métricas em mm.
\end{enumerate}

No teste da configuração atual, a chuva real tem média $4{,}447$ mm e desvio-padrão
$10{,}191$ mm. A previsão tem média $13{,}635$ mm e desvio-padrão $3{,}713$ mm. O resultado é
viés de $+9{,}188$ mm, $R^2=-0{,}803$ e MSE físico de $187{,}24$ mm$^2$. Uma climatologia
estação--mês ajustada apenas no treino obtém MSE de $103{,}29$ mm$^2$.

\section{Escopo, evidência e definição de generalização}

\subsection{Arquivos e artefatos examinados}

O caminho de execução foi seguido a partir de \path{src/run_experiment.py:120}, sem inferir o
comportamento a partir de notebooks antigos. Os principais pontos auditados foram:

\begin{itemize}
  \item configuração e orquestração: \path{src/run_experiment.py:46}--\path{src/run_experiment.py:311};
  \item dataset datado: \path{src/Data/temporal_dataset.py:129}--\path{src/Data/temporal_dataset.py:567};
  \item agregações ERA5: \path{src/Data/feature_extraction.py:173}--\path{src/Data/feature_extraction.py:257}
  e \path{src/Data/feature_extraction.py:411}--\path{src/Data/feature_extraction.py:495};
  \item topologia: \path{src/Graph/graph_related_utils.py:250}--\path{src/Graph/graph_related_utils.py:278};
  \item GLSTM: \path{src/Models/model.py:129}--\path{src/Models/model.py:306};
  \item loss e construção do modelo: \path{src/Training/experiment_runner.py:95}--\path{src/Training/experiment_runner.py:245};
  \item treino, validação e checkpoint: \path{src/Training/Training_Routines.py:40}--\path{src/Training/Training_Routines.py:350};
  \item inversão de escala e métricas físicas:
  \path{src/Evaluation/experiment_outputs.py:118-177};
  \item exportação das previsões:
  \path{src/Evaluation/experiment_outputs.py:610-655}.
\end{itemize}

Foi usada como reprodução principal a execução
\path{glstm_sweep_20260715_120845}, concluída com a configuração que está no topo do runner.
Ela produziu 576.910 pares alvo--previsão no teste. Runs anteriores foram usados somente como
ablação ou evidência histórica, e são identificados explicitamente quando aparecem.

\subsection{O que ``dados novos'' significa neste runner}

O runner mede apenas \textbf{generalização temporal para datas futuras nas mesmas 62
estações}. A função \code{chronological_split} corta o eixo do tempo, mas preserva todos os nós
nos três blocos. Portanto:

\begin{itemize}
  \item o teste é válido para novos dias nas posições já conhecidas;
  \item não existe teste de estação, região ou célula ERA5 nunca vista;
  \item o modelo possui uma adjacência $62\times 62$ e não aceita um novo nó sem reconstrução;
  \item para alegar generalização espacial seria necessário um holdout por estação/região e uma
  arquitetura indutiva para reconstruir as arestas.
\end{itemize}

Há ainda uma distinção de domínio importante: o catálogo INMET fornece nomes e coordenadas,
mas as séries de entrada e o alvo são valores ERA5 amostrados no ponto de grade mais próximo.
Em \path{src/Data/temporal_dataset.py:275}--\path{src/Data/temporal_dataset.py:294}, o próprio
canal \code{tp} do ERA5 origina $y$. Observações pluviométricas dos CSVs INMET não entram no
alvo. Assim, desempenho em uma estação real não é medido por este teste; seria uma transferência
ERA5$\rightarrow$INMET adicional.

\section{Rastreio do pipeline executado}

\subsection{Fluxo efetivo}

\begin{center}
\begin{tabularx}{0.98\textwidth}{@{}p{0.07\textwidth}Y p{0.25\textwidth}@{}}
\toprule
\textbf{Etapa} & \textbf{Operação confirmada} & \textbf{Evidência principal} \\
\midrule
1 & Lê o catálogo, filtra RS e retém nomes, latitude e longitude. &
\path{run_experiment.py:176-179}; \path{feature_extraction.py:454-467} \\
2 & Abre caches diários ERA5 e amostra a grade na coordenada mais próxima de cada estação. &
\path{temporal_dataset.py:180-192,248-376} \\
3 & Concatena 25 features e usa \code{tp} como alvo diário. &
$X=[9497,62,25]$, $y=[9497,62]$ \\
4 & Faz split cronológico no eixo bruto antes das janelas. &
\path{temporal_dataset.py:379-408} \\
5 & Cria janelas separadamente: $[D-34,\ldots,D]\mapsto[D+1,\ldots,D+5]$. &
\path{temporal_dataset.py:411-486} \\
6 & Ajusta StandardScaler de features e alvo somente nas janelas de treino. &
\path{temporal_dataset.py:508-555} \\
7 & Converte para Torch e cria DataLoaders; somente treino é embaralhado. &
\path{temporal_dataset.py:558-567}; \path{prepare_data.py:192-229} \\
8 & Constrói KNN, GLSTM, loss, Adam, scheduler e early stopping. &
\path{run_experiment.py:214-277} \\
9 & Avalia o melhor checkpoint, coleta previsões e desfaz o scaler para os CSVs/gráficos. &
\path{run_experiment.py:279-308} \\
\bottomrule
\end{tabularx}
\end{center}

\subsection{Configuração atual}

\begin{center}
\begin{tabularx}{0.94\textwidth}{@{}p{0.28\textwidth}Y@{}}
\toprule
\textbf{Componente} & \textbf{Valor efetivo} \\
\midrule
Período & 01/01/2000 a 31/12/2025; 9.497 dias \\
Nós e features & 62 nomes de estação; 25 canais ERA5 \\
Janela e horizonte & 35 dias passados; saídas D+1 a D+5 \\
Split & 60\% treino, 20\% validação, 20\% teste \\
Escala & StandardScaler global por feature e StandardScaler global do alvo \\
Grafo & KNN $k=2$; \code{LEARN_ADJ=True}; \code{LOCK_TOPOLOGY=False} \\
GLSTM & hidden 32; 1 camada; dropout 0,3; saída linear direta de 5 leads \\
Otimização & Adam; LR $10^{-3}$; weight decay $5\times10^{-3}$; clipping 1 \\
Seleção & até 100 épocas; paciência 10; menor loss ponderada de validação \\
Loss & \code{quantile_mse}; Q90; pesos inversos à frequência; máximo 20 \\
\bottomrule
\end{tabularx}
\end{center}

\section{Integridade temporal e estatística dos dados}

\subsection{Pontos que passaram na auditoria}

\callout[greenbox]{
\textbf{Não foi encontrado vazamento temporal ou desalinhamento do target.} Foram verificados
9.497 timestamps diários consecutivos, sem duplicatas, e finitude de 14.720.350 valores de $X$
e 588.814 valores de $y$. O canal histórico \code{tp} coincide com $y$ no mesmo dia, porém
somente dias até $D$ entram no input; os alvos começam em $D+1$.
}

Os intervalos brutos e as primeiras datas-alvo após o warm-up são:

\begin{center}
\begin{tabular}{@{}lccc@{}}
\toprule
\textbf{Bloco} & \textbf{Intervalo bruto} & \textbf{Primeiro alvo} & \textbf{Janelas} \\
\midrule
Treino & 01/01/2000--07/08/2015 & 05/02/2000 & 5.659 \\
Validação & 08/08/2015--18/10/2020 & 12/09/2015 & 1.860 \\
Teste & 19/10/2020--31/12/2025 & 23/11/2020 & 1.861 \\
\bottomrule
\end{tabular}
\end{center}

Os scalers e os limiares da loss são calculados somente com treino. A interseção de tempo e
estações é feita antes do split, e a conversão a tensor ocorre depois da preservação das datas em
xarray. Esses pontos são adequados.

\subsection{Distribuição do alvo e drift}

\begin{center}
\begin{tabular}{@{}lrr@{}}
\toprule
\textbf{Estatística física} & \textbf{Treino} & \textbf{Teste} \\
\midrule
Média (mm) & 4,580 & 4,447 \\
Mediana (mm) & 0,299 & 0,221 \\
Desvio-padrão (mm) & 9,745 & 10,191 \\
P90 (mm) & 15,202 & 14,474 \\
P99 (mm) & 47,018 & 50,107 \\
Máximo (mm) & 184,666 & 184,945 \\
Fração $\leq0{,}1$ mm & 42,37\% & 44,41\% \\
\bottomrule
\end{tabular}
\end{center}

O teste possui cauda levemente mais pesada, mas não há um drift grosseiro: média e escala são
próximas e a maior mudança média entre features, medida na escala do treino, foi 0,077
desvio-padrão. Além disso, o $R^2$ de treino já é negativo. Portanto, distribuição futura
diferente não é a explicação principal.

\subsection{Redundância das janelas}

Janelas consecutivas compartilham 34 de 35 dias de input (97,1\%) e 4 de 5 dias de alvo
(80\%). As 5.659 janelas de treino contêm 198.065 posições-dia de input, mas apenas 5.693 dias
distintos nessas posições. Um mesmo alvo central aparece em até cinco amostras.

Isso não cria leakage entre splits, porém cria pseudorreplicação, reduz o tamanho amostral
efetivo e torna a incerteza dos resultados muito maior do que a contagem de janelas sugere. O
scaler e os quantis também ponderam repetidamente os dias centrais. A avaliação deve usar blocos
temporais e rolling-origin, não interpretar cada janela como observação independente.

\section{Evidência quantitativa do colapso}

\subsection{Resultado da configuração atual}

Todas as métricas desta seção foram recalculadas no espaço físico a partir de
\path{test_predictions_by_lead_day.csv}. O arquivo \code{test_metrics.json} contém MSE e MAE
na escala padronizada quando \code{NORMALIZE_TARGET=True}; ele não deve ser comparado
diretamente com valores em mm.

\begin{center}
\small
\begin{tabular}{@{}crrrrrrr@{}}
\toprule
\textbf{Lead} & $\overline y$ & $\overline{\hat y}$ & $s_{\hat y}/s_y$ & \textbf{Corr.} & \textbf{Inclinação} & \textbf{Viés} & $R^2$ \\
 & \multicolumn{2}{c}{\textbf{mm}} & & & & \textbf{mm} & \\
\midrule
D+1 & 4,436 & 11,525 & 0,614 & 0,363 & 0,223 & +7,089 & -0,416 \\
D+2 & 4,445 & 13,898 & 0,328 & 0,180 & 0,059 & +9,453 & -0,850 \\
D+3 & 4,452 & 14,186 & 0,239 & 0,131 & 0,031 & +9,734 & -0,906 \\
D+4 & 4,454 & 14,135 & 0,199 & 0,118 & 0,024 & +9,681 & -0,894 \\
D+5 & 4,448 & 14,432 & 0,166 & 0,120 & 0,020 & +9,985 & -0,947 \\
\bottomrule
\end{tabular}
\end{center}

O padrão é inequívoco: a média prevista sobe para a região 14--15 mm e o desvio-padrão
desaparece conforme o lead aumenta. Em D+5, a inclinação da regressão
$\widehat y=\alpha+\beta y$ é apenas $0{,}020$; isto é, elevar a chuva real em 50 mm altera a
previsão em aproximadamente 1 mm.

No agregado dos cinco leads:

\begin{itemize}
  \item média real 4,447 mm versus média prevista 13,635 mm;
  \item desvio real 10,191 mm versus desvio previsto 3,713 mm;
  \item correlação 0,196 e inclinação 0,071;
  \item viés $+9{,}188$ mm, MAE 12,105 mm, MSE 187,241 mm$^2$ e $R^2=-0{,}803$;
  \item 44,4\% dos alvos são $\leq0{,}1$ mm, mas apenas 0,022\% das previsões;
  \item o modelo prevê mais de 10 mm em 88,5\% dos casos; a precisão desse evento é somente
  15,2\%; nenhuma previsão supera 50 mm.
\end{itemize}

\begin{figure}[H]
  \centering
  \includegraphics[width=\textwidth]{../Experiments/run_experiment/glstm_sweep_20260715_120845/prediction_overview/02_predictions_vs_actual.png}
  \caption{Run atual. Em D+5, a série prevista permanece próxima de 14 mm enquanto a chuva
  real alterna entre dias secos e picos. O scatter mostra forte compressão vertical.}
  \label{fig:current-collapse}
\end{figure}

\subsection{Baselines ajustados somente no treino}

\begin{center}
\small
\begin{tabular}{@{}lrrrrr@{}}
\toprule
\textbf{Preditor no teste} & \textbf{MSE} & \textbf{RMSE} & \textbf{MAE} & $R^2$ & \textbf{Viés} \\
 & \textbf{mm$^2$} & \textbf{mm} & \textbf{mm} & & \textbf{mm} \\
\midrule
GLSTM atual & 187,241 & 13,684 & 12,105 & -0,803 & +9,188 \\
Média global de treino & 103,884 & 10,192 & 6,139 & -0,000 & +0,135 \\
Média por estação & 103,470 & 10,172 & 6,106 & 0,004 & +0,135 \\
Climatologia estação--mês & \textbf{103,286} & \textbf{10,163} & \textbf{6,094} & \textbf{0,006} & +0,146 \\
Zero & 123,640 & 11,119 & \textbf{4,447} & -0,190 & -4,447 \\
Persistência no fim da janela & 192,750 & 13,883 & 6,910 & -0,856 & -0,019 \\
Média ponderada Q90 de treino & 221,316 & 14,877 & 13,697 & -1,131 & +10,837 \\
\bottomrule
\end{tabular}
\end{center}

A GLSTM atual tem MSE 81,3\% maior que a climatologia estação--mês. A persistência tem MSE
alto para chuva diária, mas ainda obtém MAE muito menor que a GLSTM. O baseline zero obtém MAE
de 4,447 mm porque há muitos dias secos; isso mostra por que uma única métrica não é suficiente,
mas também torna injustificável um modelo com viés positivo de quase 10 mm.

\subsection{A regressão à média sob MSE comum}

Um run histórico com MSE comum,
\path{glstm_sweep_20260713_160748}, usa outra configuração e, portanto, não é uma ablação
perfeitamente controlada. Ele serve para mostrar o mecanismo clássico: quando a perda volta a
ser MSE e o sinal desaparece, a constante deixa de ser 15 mm e volta à média de aproximadamente
4--5 mm.

\begin{center}
\small
\begin{tabular}{@{}crrrrr@{}}
\toprule
\textbf{Lead} & $\overline y$ & $\overline{\hat y}$ & $s_{\hat y}/s_y$ & \textbf{Corr.} & $R^2$ \\
\midrule
D+1 & 4,806 & 4,321 & 0,394 & 0,429 & 0,181 \\
D+2 & 4,817 & 4,053 & 0,151 & 0,177 & 0,026 \\
D+3 & 4,824 & 4,027 & 0,102 & 0,112 & 0,007 \\
D+4 & 4,828 & 4,108 & 0,076 & 0,092 & 0,004 \\
D+5 & 4,831 & 4,154 & 0,065 & 0,094 & 0,004 \\
\bottomrule
\end{tabular}
\end{center}

\begin{figure}[H]
  \centering
  \includegraphics[width=\textwidth]{../Experiments/run_experiment/glstm_sweep_20260713_160748/prediction_overview/02_predictions_vs_actual.png}
  \caption{Run histórico com MSE. Em D+5, a saída laranja converge para a média comum e quase
  não varia. A loss atual muda a altura da linha, mas não resolve a falta de dispersão.}
  \label{fig:mse-collapse}
\end{figure}

\section{Por que a loss produz uma média ponderada}

\subsection{Derivação}

Para uma constante $c$ e MSE ponderada,

\[
L(c)=\frac{\sum_i w_i(c-y_i)^2}{\sum_i w_i},
\qquad
\frac{\partial L}{\partial c}=0
\quad\Longrightarrow\quad
c^*=\frac{\sum_i w_i y_i}{\sum_i w_i}.
\]

Na forma condicional, o regressor ótimo é

\[
f^*(x)=\frac{\mathbb{E}[w(Y)Y\mid X=x]}{\mathbb{E}[w(Y)\mid X=x]}.
\]

Se $X$ informa pouco sobre o futuro, $f^*(x)$ varia pouco e se aproxima do centroide global.
Em MSE comum, $w=1$ e esse centroide é a média. Na loss atual, os pesos são construídos em
\path{src/Training/experiment_runner.py:169}--\path{src/Training/experiment_runner.py:235}:

\begin{itemize}
  \item 90\% inferiores: peso $0{,}199995$, média $1{,}902$ mm;
  \item 10\% superiores ao Q90: peso $1{,}800005$, média $28{,}681$ mm.
\end{itemize}

Cada extremo recebe cerca de nove vezes o peso de um caso comum. Como há nove vezes mais casos
comuns, os dois grupos recebem aproximadamente a mesma massa total. O ótimo constante torna-se

\[
c^*\approx\frac{1{,}902+28{,}681}{2}=15{,}291\ \text{mm}.
\]

É exatamente a região para a qual D+2--D+5 convergem. O nome
\code{quantile_mse} pode induzir a interpretação errada: ele \textbf{não} estima o quantil 0,9.
Para isso seria necessária pinball loss. Ele estima uma média condicional reponderada.

\subsection{A seleção do checkpoint reforça o desalinhamento}

O early stopping e o scheduler usam a mesma loss ponderada. No run atual:

\begin{center}
\begin{tabular}{@{}lrrrr@{}}
\toprule
\textbf{Estado} & \textbf{Época} & \textbf{Val loss} & \textbf{Val MSE pad.} & \textbf{Val $R^2$} \\
\midrule
Primeira época & 1 & 2,7769 & 1,9263 & -0,898 \\
Melhor MSE/$R^2$ ordinário & 6 & 2,7365 & \textbf{1,7674} & \textbf{-0,739} \\
Checkpoint pela loss configurada & 13 & \textbf{2,7121} & 1,9842 & -0,952 \\
Última época registrada & 24 & 2,7204 & 1,9869 & -0,954 \\
\bottomrule
\end{tabular}
\end{center}

A loss ponderada melhora, enquanto o $R^2$ físico piora. O otimizador está fazendo o que foi
pedido. O problema é que o critério de seleção não corresponde ao comportamento que se deseja
avaliar.

\begin{figure}[H]
  \centering
  \includegraphics[width=0.92\textwidth]{../Experiments/run_experiment/glstm_sweep_20260715_120845/r2_curve.png}
  \caption{$R^2$ ordinário negativo em treino e validação no run atual. Não há o padrão
  principal de overfit clássico; o objetivo monitorado é outro.}
\end{figure}

\section{O grafo aprendível está promovendo oversmoothing}

\subsection{Mecanismo no código}

Em \path{src/Models/model.py:139}--\path{src/Models/model.py:187}, uma não-aresta é
representada por probabilidade $10^{-4}$, isto é, logit aproximadamente $-9{,}21$. Com
\code{LOCK_TOPOLOGY=False}, todas as 3.624 não-arestas entram no vetor de parâmetros.

O treino separa \code{a_logits} apenas para permitir outro learning rate, mas cria o Adam com o
mesmo \code{weight_decay} para todos os grupos
(\path{src/Training/Training_Routines.py:126}--\path{src/Training/Training_Routines.py:145}).
L2 leva logits para zero. Para uma aresta probabilística, isso não é esparsidade:

\[
a_{ij}\rightarrow 0
\quad\Longrightarrow\quad
\sigma(a_{ij})\rightarrow 0{,}5.
\]

Depois de cada passo, os logits também são truncados em $[-6,6]$
(\path{src/Training/Training_Routines.py:226}--\path{src/Training/Training_Routines.py:229}).
O primeiro clamp sozinho eleva uma não-aresta de $10^{-4}$ para
$\sigma(-6)=0{,}00247$, cerca de 25 vezes.

\subsection{Evidência no checkpoint}

\begin{center}
\begin{tabular}{@{}lrr@{}}
\toprule
\textbf{Componente da adjacência} & \textbf{Inicial} & \textbf{Checkpoint} \\
\midrule
Probabilidade média nas 158 arestas dirigidas KNN & 0,9500 & 0,5433 \\
Probabilidade média nas 3.624 não-arestas & 0,0001 & 0,09049 \\
Massa normalizada em self-loops & --- & 13,05\% \\
Massa normalizada nas arestas KNN & --- & 17,87\% \\
Massa normalizada nas não-arestas & --- & \textbf{69,04\%} \\
Segundo autovalor em módulo $|\lambda_2|$ & --- & 0,2823 \\
Proxy $|\lambda_2|^{35}$ & --- & $\approx6\times10^{-20}$ \\
\bottomrule
\end{tabular}
\end{center}

O proxy espectral não é uma igualdade para a GLSTM, porque gates e transformações não lineares
intervêm. Ele mostra, contudo, quão rapidamente a propagação espacial pura apagaria componentes
não constantes. Isso é relevante porque \code{H_prev @ A} e \code{C_prev @ A} são aplicados em
cada passo da célula (\path{src/Models/model.py:197}--\path{src/Models/model.py:215}) e o loop
tem 35 passos (\path{src/Models/model.py:277}--\path{src/Models/model.py:280}).

A dispersão espacial média no alvo é 4,13 mm; ela cai para 1,25 mm no modelo com grafo
aprendível e 1,56 mm no par com grafo fixo. Em uma comparação quase controlada,
\path{glstm_sweep_20260715_103451} (grafo fixo) obteve $R^2=-0{,}747$, enquanto o run atual
aprendível obteve $-0{,}803$. Logo, o grafo não é a causa única, mas piora a homogeneização.

\begin{figure}[H]
  \centering
  \begin{minipage}[t]{0.49\textwidth}
    \centering
    \includegraphics[width=\textwidth]{../Experiments/run_experiment/glstm_sweep_20260715_120845/initial_topology.png}
    \caption*{Inicial: KNN forte e não-arestas próximas de zero.}
  \end{minipage}\hfill
  \begin{minipage}[t]{0.49\textwidth}
    \centering
    \includegraphics[width=\textwidth]{../Experiments/run_experiment/glstm_sweep_20260715_120845/topology.png}
    \caption*{Final: arestas originais enfraquecidas e fundo denso de baixa intensidade.}
  \end{minipage}
  \caption{Adjacência não normalizada antes e depois do treinamento. A escala visual escura
  torna os pesos 0,09 discretos, mas a quantidade de não-arestas faz sua massa dominar.}
\end{figure}

\section{Limitação informacional e arquitetural}

\subsection{O input não é uma previsão meteorológica de D+1--D+5}

As 25 features são agregados diários históricos de precipitação, temperatura/ponto de orvalho,
umidade específica e velocidade vertical. \code{INCLUDE_WIND=False}. Não há:

\begin{itemize}
  \item campos NWP previstos para cada lead;
  \item seno/cosseno de dia do ano ou mês;
  \item latitude, longitude, altitude ou embedding persistente da estação;
  \item índices de circulação, pressão à superfície, CAPE ou umidade integrada configurados no
  run atual;
  \item direção espacial da advecção; o grafo é estático e simétrico.
\end{itemize}

O catálogo contém altitudes de 0 a 1.228,59 m, porém
\path{src/Data/feature_extraction.py:454}--\path{src/Data/feature_extraction.py:467} retém somente
latitude e longitude, e estas nem sequer entram como feature do modelo. O tempo permanece em
xarray como coordenada, mas é descartado na fronteira Torch. Assim, a rede precisa inferir
sazonalidade e orografia apenas pelas séries passadas.

Uma verificação univariada confirma o rápido desaparecimento do sinal: a persistência da última
chuva observada tem correlação 0,220 com D+1 e 0,042 com D+5. A maior correlação absoluta entre
qualquer feature no fim da janela e o alvo foi 0,288 em D+1 e 0,049 em D+5. Relações não
lineares podem existir, mas esses valores explicam por que um regressor determinístico se apoia
no centroide para leads longos.

\subsection{Head direto e gargalo comum}

Em \path{src/Models/model.py:253}--\path{src/Models/model.py:281}, o último estado oculto de
cada nó passa por um MLP que produz os cinco leads de uma vez. Não há decoder temporal,
embedding de lead, mecanismo autoregressivo ou covariável futura. A mesma representação deve
servir simultaneamente D+1 e D+5.

Isso produz a sequência observada: correlação 0,363 e razão de desvios 0,614 em D+1, mas
correlação 0,120 e razão 0,166 em D+5. Mesmo sob MSE simples, D+5 conserva somente 6,5\% do
desvio real. Aumentar hidden de 32 para 256, mantendo uma configuração próxima, não resolveu
($R^2$ passou de $-0{,}747$ para $-0{,}776$). Capacidade bruta não substitui informação e um
objetivo adequado.

\subsection{A natureza de uma previsão pontual}

Mesmo um regressor MSE correto estima $\mathbb{E}[Y\mid X]$, não uma realização aleatória de
$Y$. Portanto, sua série deve ser mais suave que a observada. O problema aqui não é a mera
existência de suavização; é sua magnitude e sua falta de skill: inclinações próximas de zero,
$R^2$ negativo e desempenho inferior a climatologias de treino.

Para precipitação zero-inflada e de cauda longa, uma previsão distribucional ou em duas partes
é mais informativa que forçar um único valor a representar simultaneamente ocorrência, volume e
extremos.

\section{Problemas secundários e observabilidade}

\begin{longtable}{@{}p{0.27\textwidth}p{0.67\textwidth}@{}}
\toprule
\textbf{Achado} & \textbf{Impacto} \\
\midrule
\endfirsthead
\toprule
\textbf{Achado} & \textbf{Impacto} \\
\midrule
\endhead
KNN em graus & \path{graph_related_utils.py:250-278} usa distância euclidiana em
[longitude, latitude], sem Haversine nem peso por distância. Com $k=2$, 7 de 62 nós mudam de
vizinho em relação ao KNN geodésico; as arestas chegam a 150,9 km. \\
\addlinespace
Nós duplicados & As 62 coordenadas de estação mapeiam para apenas 60 células ERA5 de 0,25 grau.
BAGE/BAGE--CENTRO e CHARQUEADAS/PARQUE ELDORADO têm séries idênticas, embora sejam nós
separados. \\
\addlinespace
Scaler global do alvo & \path{temporal_dataset.py:529-550} usa \code{reshape(-1,1)} para todas
as estações e leads. Se o sinal se perde, zero padronizado volta diretamente à média global, não
a uma climatologia específica do nó/estação. \\
\addlinespace
Unidades stale & A precipitação é multiplicada por 1.000 em
\path{feature_extraction.py:438-443}, mas o Zarr diário ainda declara unidade ``m''; temperatura
é convertida a graus Celsius, enquanto o cache ainda pode declarar ``K''. Os valores usados no
run são coerentes com mm/°C, mas o metadado cria risco de dupla conversão futura. \\
\addlinespace
Métricas em escalas diferentes & \code{test_metrics.json} mede loss/MSE/MAE no alvo
padronizado; os CSVs posteriores são inversamente transformados para mm. O nome das métricas
não carrega a unidade. \\
\addlinespace
MAPE sem interpretação & Com muitos zeros, MAPE físico já é instável. No treino atual ele é
calculado sobre o alvo padronizado; zero significa ``igual à média do treino'', não ``sem chuva''.
O valor 476,97\% não deve orientar seleção. \\
\addlinespace
Média por batch & Loss, MSE e MAE são promediadas por número de batches, não pelo número de
elementos. O último batch menor recebe peso excessivo. O $R^2$ global usa acumuladores corretos. \\
\addlinespace
Resumo versus checkpoint & \code{run_summary.json} guarda várias métricas da última época,
mas \code{model_state_dict.pt} guarda o melhor checkpoint pela loss. As duas linhas não
descrevem necessariamente o mesmo modelo. \\
\addlinespace
Datas nulas no resumo & Os DataLoaders usam \code{TensorDataset}; por isso
\code{train_start_date}, \code{train_end_date} e equivalentes saem nulos, embora as datas ainda
existam em \code{dataset_contract.json}. \\
\addlinespace
Sem diagnóstico de clipping & O retorno de \code{clip_grad_norm_} é descartado. Não se sabe em
quantos batches o limite 1 está efetivamente reduzindo o gradiente. \\
\bottomrule
\end{longtable}

\section{Hipóteses descartadas ou rebaixadas}

\begin{itemize}
  \item \textbf{Leakage/target deslocado:} não confirmado; as datas e shapes estão corretos e os
  alvos de cada split são disjuntos.
  \item \textbf{NaN/Inf:} não confirmado nos dados atuais.
  \item \textbf{Drift como causa dominante:} rebaixado; treino e teste têm estatísticas
  próximas, e o treino também apresenta $R^2$ negativo.
  \item \textbf{Overfit clássico do default:} rebaixado; loss de treino e validação são próximas.
  Há runs históricos maiores com gap treino--validação, mas esse não é o padrão principal do run
  atual.
  \item \textbf{Normalização do target como causa isolada:} rebaixada; ablações com e sem
  normalização permanecem ruins. O problema é o objetivo, o sinal e a arquitetura espacial.
  \item \textbf{Poucos parâmetros apenas:} rebaixado; hidden 256 não trouxe ganho. Um modelo
  maior pode memorizar treino sem resolver a transferência temporal.
\end{itemize}

\section{Plano recomendado}

\subsection{Prioridade P0 --- tornar a conclusão mensurável}

Antes de nova arquitetura, o runner deve produzir em todo run:

\begin{enumerate}
  \item baselines ajustados somente no treino: zero, média global, média por estação,
  climatologia estação--mês e persistência no fim da janela;
  \item métricas físicas por lead, estação, estação do ano e regime de chuva: MSE, RMSE, MAE,
  $R^2$, viés, correlação, inclinação e $s_{\hat y}/s_y$;
  \item classificação para limiares como 0,1, 1, 10 mm, Q90 e Q95: precisão, recall, FAR, CSI e
  Brier score quando houver probabilidade;
  \item datas exatas dos splits, scaler completo, unidade de cada artefato, melhor época em cada
  métrica e norma pré-clipping;
  \item gráficos da massa em arestas/não-arestas, grau efetivo, $|\lambda_2|$ e variância espacial
  por época.
\end{enumerate}

\subsection{Prioridade P1 --- separar objetivo, horizonte e grafo}

Executar uma matriz controlada, mesma seed, mesmo split e mesmos dados:

\begin{center}
\small
\begin{tabularx}{\textwidth}{@{}p{0.09\textwidth}p{0.19\textwidth}p{0.17\textwidth}Y@{}}
\toprule
\textbf{Teste} & \textbf{Horizonte} & \textbf{Grafo} & \textbf{Objetivo} \\
\midrule
A0 & --- & --- & Climatologia estação--mês \\
A1 & D+1 & Sem grafo & MSE comum \\
A2 & D+1 & KNN Haversine fixo & MSE comum \\
A3 & D+1 & Topologia KNN travada & MSE, adjacência aprendível sem decay nos logits \\
A4 & D+1--D+5 & Melhor opção A1--A3 & Um modelo separado por lead ou decoder com embedding de lead \\
A5 & D+1--D+5 & Melhor grafo & Pinball verdadeira para quantis explícitos \\
A6 & D+1--D+5 & Melhor grafo & Ocorrência + quantidade positiva/distribuição \\
\bottomrule
\end{tabularx}
\end{center}

Não avançar para D+5 antes de obter skill reprodutível em D+1. Um teste adicional deve fazer a
rede memorizar deliberadamente 16--32 janelas, sem dropout, weight decay, AMP ou grafo
aprendível. Se ela não levar a loss quase a zero, existe um problema de capacidade/gradiente
separado da generalização.

\subsection{Prioridade P1 --- corrigir o grafo aprendível}

\begin{itemize}
  \item usar \code{LEARN_ADJ=False} como baseline obrigatório;
  \item quando aprender pesos, começar com \code{LOCK_TOPOLOGY=True};
  \item definir \code{weight_decay=0} especificamente para \code{a_logits};
  \item substituir L2 nos logits por penalidade de esparsidade sobre probabilidades off-diagonal,
  top-$k$ ou orçamento de massa explícito;
  \item não usar clamp incompatível com a probabilidade inicial $10^{-4}$; ou inicializar de modo
  coerente com o intervalo permitido;
  \item comparar KNN Haversine ponderado por distância e grafo dirigido/dinâmico informado por
  vento;
  \item testar mistura espacial residual uma vez por bloco, em vez de difusão total do estado de
  célula em todos os passos.
\end{itemize}

Um teste unitário conceitual deve executar um passo sem gradiente informativo e confirmar que
uma não-aresta não cresce apenas por weight decay/clamp. O comportamento atual falha nesse
critério.

\subsection{Prioridade P2 --- fornecer informação realmente preditiva}

\begin{itemize}
  \item incluir seno/cosseno de dia do ano e, se apropriado, hora/estação;
  \item incluir latitude, longitude, altitude e/ou embedding de nó;
  \item habilitar e validar vento, preferencialmente de modo direcional;
  \item para previsão D+1--D+5, fornecer covariáveis NWP válidas em cada lead, não somente
  observações históricas;
  \item considerar previsão residual sobre climatologia ou NWP, preservando um caminho de skip;
  \item se o alvo final for estação INMET, treinar e avaliar contra observação INMET, com ERA5
  como covariável, não como substituto silencioso do alvo.
\end{itemize}

\subsection{Prioridade P2 --- objetivo adequado à precipitação}

Há três perguntas distintas e elas não devem ser misturadas em uma única MSE ponderada:

\begin{enumerate}
  \item \textbf{Vai chover?} Classificação/probabilidade de ocorrência.
  \item \textbf{Quanto, dado que chove?} Distribuição positiva, por exemplo Gamma/lognormal, ou
  uma família como Tweedie quando apropriada.
  \item \textbf{Qual o risco de extremo?} Quantis com pinball loss ou probabilidade de exceder
  limiares, avaliada com métricas probabilísticas.
\end{enumerate}

\code{Softplus} não deve ser aplicado diretamente a um alvo padronizado que possui valores
negativos. A transformação e a distribuição precisam ser definidas de forma consistente. Uma
saída determinística pode permanecer como baseline, mas deve ser comparada à climatologia e não
ser tratada como amostra realista da chuva.

\subsection{Critérios de aceitação}

Uma modificação só deve ser considerada melhoria se:

\begin{itemize}
  \item superar climatologia estação--mês no mesmo bloco temporal, por lead;
  \item manter viés e razão de desvios reportados, evitando ganho aparente por previsão constante;
  \item repetir o ganho em múltiplas seeds e janelas rolling-origin;
  \item para a GNN, superar o mesmo modelo sem grafo sob idênticas condições;
  \item se a alegação incluir novas estações, superar um holdout espacial explícito;
  \item preservar esparsidade/massa do grafo dentro de limites definidos antes do treino;
  \item reportar separadamente skill em seco, chuva moderada e extremos.
\end{itemize}

\section{Resposta direta}

O modelo converge para a média porque uma previsão pontual por erro quadrático escolhe o
centro condicional quando não consegue distinguir os futuros possíveis. No GLSTM atual, essa
tendência é amplificada por quatro fatos concretos:

\begin{enumerate}
  \item o histórico contém pouco sinal para chuva exata em D+2--D+5 e não há preditores futuros;
  \item os cinco leads saem do mesmo estado oculto final, cuja variabilidade diminui com o lead;
  \item o grafo aprendível torna-se quase global e mistura os nós repetidamente;
  \item a loss atual troca a média comum de 4,58 mm por uma média ponderada de 15,29 mm.
\end{enumerate}

Por isso, com MSE comum vê-se uma linha próxima de 4--5 mm; com
\code{quantile_mse}, vê-se uma linha próxima de 14--15 mm. O fenômeno é o mesmo: pouca
dependência de $X$ e grande dependência do centroide imposto pela função objetivo.

\callout[amberbox]{
\textbf{Síntese final.} O primeiro passo não é aumentar a rede. É alinhar objetivo e métrica,
desativar a adjacência livre, provar skill em D+1 contra climatologia e adicionar informação
futura/estática. Só depois faz sentido reintroduzir um grafo aprendível e o horizonte completo.
}

\appendix
\section{Trilha de evidências reproduzíveis}

\begin{longtable}{@{}p{0.33\textwidth}p{0.61\textwidth}@{}}
\toprule
\textbf{Artefato} & \textbf{Uso nesta auditoria} \\
\midrule
\endfirsthead
\toprule
\textbf{Artefato} & \textbf{Uso nesta auditoria} \\
\midrule
\endhead
\path{.../120845/config.json} & Configuração efetiva do run atual. \\
\path{.../120845/dataset_contract.json} & Shapes, features, estações e cobertura temporal. \\
\path{.../120845/run_summary.json} & Pesos/limiar da loss, épocas, loss e $R^2$ registrados. \\
\path{.../120845/hist.pt} & Localização das melhores épocas por loss e por MSE/$R^2$. \\
\path{.../120845/model_state_dict.pt} & Probabilidades e espectro da adjacência do checkpoint. \\
\path{.../120845/test_predictions_by_lead_day.csv} & Distribuições, correlação, inclinação, viés e métricas físicas. \\
\path{.../120845/forecast_horizon_diagnostics/test_prediction_metrics_by_lead_day.csv} & Conferência das métricas físicas por lead. \\
\path{.../160748/test_predictions_by_lead_day.csv} & Evidência histórica do colapso à média sob MSE comum. \\
\path{Datasets/processed/daily/precipitation_daily.zarr} & Reconstrução de distribuição de treino e baselines sem usar o teste para ajuste. \\
\bottomrule
\end{longtable}

\section{Limitações da auditoria}

\begin{itemize}
  \item As comparações identificadas como históricas usam hiperparâmetros diferentes e não são
  apresentadas como efeito causal isolado.
  \item Correlações univariadas não excluem relações não lineares; elas medem a dificuldade e a
  perda progressiva de sinal, não uma impossibilidade matemática.
  \item Os totais diários ERA5 não foram confrontados com uma fonte meteorológica independente
  nem com pluviômetros INMET nesta auditoria.
  \item Não foi implementada uma correção funcional. O documento registra o diagnóstico e uma
  sequência de experimentos capaz de separar as causas.
\end{itemize}

\vfill
\begin{center}
\color{muted}\small
Documento gerado a partir do código e dos artefatos locais do repositório em 15/07/2026.
\end{center}

\end{document}
